\begin{enumerate}[start=2]
\item Comme $a(t)=v'(t)$ et que $v(t)=2t$ sur l'intervalle [1 ; 4], on obtient :$a(t)=v'(t)=2$.

L'accélération est donc constante et vaut $2$~m$\cdot$s$^{-2}$ sur cet intervalle.

\item La distance parcourue correspond à l'aire sous la courbe de la vitesse. Le carré hachuré a pour dimensions $5$~s et $5$~m$\cdot$s$^{-1}$, donc son aire représente : $5\times 5=25~\text{m}$.

En comptant les carreaux sous la courbe, on obtient environ $23$ carrés de ce type.

La distance totale parcourue est donc d'environ $23\times 25=575$~m.

\item La vitesse moyenne du véhicule sur l'intervalle [0 ; 50] est d'environ :
\[v_{moy} \approx \dfrac{575}{50} \times \dfrac{60 \times 60}{\np{1000}}=41,4~\text{km}\cdot\text{h}^{-1}.\]
\end{enumerate}

\bigskip
